\clearpage\gdef\curyear{2016--2026 \textperiodcentered{} Summary}
\begin{center}{\pgfornament[width=5cm,color=gilt]{87}}\par
{\fontsize{20}{24}\selectfont\color{navy}\scshape KCET 2016--2026 \textperiodcentered{} Summary}\par
{\color{gilt}\pgfornament[width=5cm,color=gilt,symmetry=h]{87}}\end{center}\label{psummary}
\vspace{6pt}
{\small Every question of the eleven papers was checked one by one against the four 2026-27 NCERT Physics books (Class XI Part 1 \& 2, Class XII Part 1 \& 2). \textsc{Direct}: the answer is stated in the textbook (text, summary, table or figure), or the question is an NCERT solved Example or Exercise (same or changed numbers), or it needs only one printed formula. \textsc{Applied}: the concept is in the textbook but the question needs several steps or a new situation. \textsc{Deleted portion}: the topic was removed from the rationalised syllabus; such questions are not analysed. Percentages are of the analysed questions.}\par\medskip
\begin{center}\small
\begin{tabular}{@{}crrrrrrrr@{}}\toprule
Year & \multicolumn{4}{c}{Direct} & Applied & Analysed & Deleted & Direct\\
\cmidrule(lr){2-5}
 & Total & Text & Example & Exercise & & & portion & \%\\\midrule

2026 & \textbf{37} & 32 & 3 & 2 & 23 & 60 & 0 & 62\%\\
2025 & \textbf{38} & 34 & 3 & 1 & 22 & 60 & 0 & 63\%\\
2024 & \textbf{24} & 21 & 2 & 1 & 28 & 52 & 8 & 46\%\\
2023 & \textbf{20} & 14 & 2 & 4 & 24 & 44 & 16 & 45\%\\
2022 & \textbf{34} & 27 & 2 & 5 & 23 & 57 & 3 & 60\%\\
2021 & \textbf{26} & 20 & 3 & 3 & 26 & 52 & 8 & 50\%\\
2020 & \textbf{25} & 18 & 4 & 3 & 24 & 49 & 11 & 51\%\\
2019 & \textbf{20} & 18 & 1 & 1 & 28 & 48 & 12 & 42\%\\
2018 & \textbf{32} & 25 & 4 & 3 & 18 & 50 & 10 & 64\%\\
2017 & \textbf{35} & 26 & 7 & 2 & 16 & 51 & 9 & 69\%\\
2016 & \textbf{27} & 22 & 4 & 1 & 21 & 48 & 12 & 56\%\\
\midrule\textbf{Total} & \textbf{318} & 257 & 35 & 26 & 253 & 571 & \textbf{89} & 56\%\\\bottomrule
\end{tabular}\end{center}\medskip
{\small Of the 660 questions, \textbf{571} are within the present syllabus. Of these, \textbf{318} (56\%) are direct from the NCERT textbook and 253 (44\%) are applications of NCERT concepts. \textbf{89} questions (13\% of all 660) are from deleted portions; most of them are in the papers of 2016--2024, which were set before the rationalised syllabus was fully followed. KCET 2025 and 2026 have no question from a deleted portion.}\par\medskip
{\small\textbf{\color{crimson}Deleted portions asked, 2016--2026} (number of questions):\par\smallskip
\begin{multicols}{2}\raggedright

Digital electronics (logic gates) (8);\\
Earth's magnetism (incl. angle of dip) (7);\\
Alpha- and beta-decay (6);\\
Radioactive decay law (half-life, mean life, activity) (6);\\
Cyclotron (5);\\
Doppler effect (5);\\
Transistor (incl. amplifier) (5);\\
Communication systems (incl. antenna) (4);\\
Resistor colour code (4);\\
Curie's law (3);\\
LC oscillations (3);\\
Hysteresis and coercivity (2);\\
Oscillation period of a magnetic needle (2);\\
Perpendicular axis theorem (2);\\
Potentiometer (2);\\
Radiation pressure of EM waves (2);\\
Rolling motion (2);\\
Sharpness of resonance (Q-factor) (2);\\
Toroid (2);\\
Zener diode (2);\\
Accuracy and precision of measurements (1);\\
Brewster's law (1);\\
Charging by induction (1);\\
Coefficient of restitution (1);\\
Curie--Weiss law (1);\\
Damped oscillations (1);\\
Eddy currents (1);\\
Fresnel distance (1);\\
Light-emitting diode (1);\\
Parallel axis theorem (1);\\
Photodiode (1);\\
Resolving power of the eye (1);\\
Scattering of light (1);\\
Velocity selector (1);\\
Voltage regulation (1).
\end{multicols}}

\clearpage
\begin{center}{\pgfornament[width=5cm,color=gilt]{87}}\par
{\fontsize{20}{24}\selectfont\color{navy}\scshape The 15 Most Repeated Topics}\par
{\color{gilt}\pgfornament[width=5cm,color=gilt,symmetry=h]{87}}\end{center}\label{ptopics}
\vspace{4pt}
{\small Every in-syllabus question of 2016--2026 was given one topic label while the answer keys were solved. The table counts the questions on each topic over the eleven papers (deleted portions left out; ties are ordered by the number of years in which the topic appeared). Years in bold are the years the topic was asked.}\par\medskip
{\small\renewcommand{\arraystretch}{1.3}
\begin{tabularx}{\linewidth}{@{}r@{\hspace{6pt}}>{\raggedright\arraybackslash}Xlcc@{\hspace{6pt}}l@{}}\toprule
\# & Topic & NCERT chapter & Qs & Years & 2016 $\rightarrow$ 2026\\\midrule

1 & Bohr model of hydrogen atom & Atoms (XII) & \textbf{20} & 11 & {\scriptsize \textbf{16} \textbf{17} \textbf{18} \textbf{19} \textbf{20} \textbf{21} \textbf{22} \textbf{23} \textbf{24} \textbf{25} \textbf{26}}\\
2 & AC circuits (reactance, LCR, phase) & Alternating Current (XII) & \textbf{16} & 11 & {\scriptsize \textbf{16} \textbf{17} \textbf{18} \textbf{19} \textbf{20} \textbf{21} \textbf{22} \textbf{23} \textbf{24} \textbf{25} \textbf{26}}\\
3 & Biot-Savart law and magnetic fields & Moving Charges \& Magnetism (XII) & \textbf{16} & 9 & {\scriptsize \textbf{16} \textbf{17} \textbf{18} \textbf{19} \textbf{20} {\color{gray!45}21} \textbf{22} {\color{gray!45}23} \textbf{24} \textbf{25} \textbf{26}}\\
4 & Gauss's law and electric flux & Electric Charges \& Fields (XII) & \textbf{16} & 8 & {\scriptsize \textbf{16} \textbf{17} {\color{gray!45}18} {\color{gray!45}19} \textbf{20} \textbf{21} {\color{gray!45}22} \textbf{23} \textbf{24} \textbf{25} \textbf{26}}\\
5 & Photoelectric effect & Dual Nature (XII) & \textbf{15} & 11 & {\scriptsize \textbf{16} \textbf{17} \textbf{18} \textbf{19} \textbf{20} \textbf{21} \textbf{22} \textbf{23} \textbf{24} \textbf{25} \textbf{26}}\\
6 & Current, drift velocity and mobility & Current Electricity (XII) & \textbf{14} & 9 & {\scriptsize \textbf{16} \textbf{17} {\color{gray!45}18} \textbf{19} {\color{gray!45}20} \textbf{21} \textbf{22} \textbf{23} \textbf{24} \textbf{25} \textbf{26}}\\
7 & Lenses and lens maker's formula & Ray Optics (XII) & \textbf{14} & 8 & {\scriptsize \textbf{16} {\color{gray!45}17} \textbf{18} {\color{gray!45}19} \textbf{20} \textbf{21} \textbf{22} \textbf{23} \textbf{24} \textbf{25} {\color{gray!45}26}}\\
8 & Kinematics in one dimension & Motion in a Straight Line (XI) & \textbf{12} & 10 & {\scriptsize \textbf{16} \textbf{17} \textbf{18} \textbf{19} {\color{gray!45}20} \textbf{21} \textbf{22} \textbf{23} \textbf{24} \textbf{25} \textbf{26}}\\
9 & Capacitors (energy, dielectrics) & Electrostatic Potential \& Capacitance (XII) & \textbf{12} & 9 & {\scriptsize \textbf{16} \textbf{17} \textbf{18} \textbf{19} {\color{gray!45}20} \textbf{21} \textbf{22} \textbf{23} \textbf{24} {\color{gray!45}25} \textbf{26}}\\
10 & Kirchhoff's rules and DC circuits & Current Electricity (XII) & \textbf{12} & 8 & {\scriptsize \textbf{16} \textbf{17} \textbf{18} \textbf{19} \textbf{20} {\color{gray!45}21} {\color{gray!45}22} \textbf{23} \textbf{24} {\color{gray!45}25} \textbf{26}}\\
11 & Electric potential and potential energy & Electrostatic Potential \& Capacitance (XII) & \textbf{11} & 10 & {\scriptsize \textbf{16} \textbf{17} \textbf{18} \textbf{19} \textbf{20} \textbf{21} \textbf{22} \textbf{23} {\color{gray!45}24} \textbf{25} \textbf{26}}\\
12 & AC power and rms values & Alternating Current (XII) & \textbf{11} & 6 & {\scriptsize \textbf{16} {\color{gray!45}17} \textbf{18} \textbf{19} {\color{gray!45}20} {\color{gray!45}21} \textbf{22} {\color{gray!45}23} {\color{gray!45}24} \textbf{25} \textbf{26}}\\
13 & Elasticity (Young's modulus) & Mech.\ Prop.\ of Solids (XI) & \textbf{10} & 10 & {\scriptsize \textbf{16} \textbf{17} \textbf{18} \textbf{19} \textbf{20} {\color{gray!45}21} \textbf{22} \textbf{23} \textbf{24} \textbf{25} \textbf{26}}\\
14 & Electric dipole & Electric Charges \& Fields (XII) & \textbf{10} & 9 & {\scriptsize \textbf{16} \textbf{17} {\color{gray!45}18} \textbf{19} \textbf{20} \textbf{21} \textbf{22} \textbf{23} {\color{gray!45}24} \textbf{25} \textbf{26}}\\
15 & Young's double-slit interference & Wave Optics (XII) & \textbf{10} & 9 & {\scriptsize \textbf{16} \textbf{17} \textbf{18} \textbf{19} \textbf{20} \textbf{21} \textbf{22} \textbf{23} {\color{gray!45}24} {\color{gray!45}25} \textbf{26}}\\
\bottomrule\end{tabularx}}\par\medskip
{\small These 15 topics carry \textbf{199} of the 571 in-syllabus questions (35\%), about 18 questions in every paper. Only Kinematics in one dimension and Elasticity (Young's modulus) are from Class XI; the rest are Class XII.}\par\medskip
{\small\textbf{\color{crimson}Questions asked again, almost word for word} (in the present syllabus)\par\smallskip
\begin{tabular}{@{}ll@{}}
2019 Q20 $=$ 2020 Q48 & two protons 10\,nm apart: nuclear force vs Coulomb force\\
2019 Q4 $=$ 2020 Q30 & permanent magnet at room temperature: domains partly aligned\\
2019 Q17 $\approx$ 2020 Q45 & frequency (2019) and period (2020) of the electron in the $n$th Bohr orbit\\
2018 Q37 $\approx$ 2026 Q44 & current in a 100\,W bulb on 220\,V, 50\,Hz AC\\
\end{tabular}}\par

\clearpage
\begin{center}{\pgfornament[width=5cm,color=gilt]{87}}\par
{\fontsize{20}{24}\selectfont\color{navy}\scshape How to Prepare for KCET Physics}\par
{\color{gilt}\pgfornament[width=5cm,color=gilt,symmetry=h]{87}}\end{center}\label{pplan}
\vspace{2pt}
{\small Every point below comes from the question-by-question analysis of the eleven papers in this book.}\par\medskip
\begin{multicols}{2}
\textbf{\color{crimson}1. What the papers show}\par\smallskip
\begin{itemize}[leftmargin=1.1em,itemsep=2pt,topsep=0pt]
\item \textbf{56\%} of the in-syllabus questions (318 of 571) are direct from the NCERT textbook; the rest (253) apply NCERT concepts. Only 42 applied questions closely resemble a NEET/JEE/AIEEE question.
\item 35 questions are NCERT solved Examples and 26 are NCERT Exercises, often with the same numbers.
\item Class XII gives \textbf{71\%} of the questions, Class XI 29\%.
\item Heaviest chapters (questions in 11 years): Current Electricity (61), Electrostatic Potential \& Capacitance (43), Moving Charges \& Magnetism (42), Ray Optics (41), Electric Charges \& Fields (34), Alternating Current (31).
\item Difficulty: 360 easy, 202 moderate, only 9 hard. KCET rewards speed and accuracy more than tricks.
\item 89 old questions are on deleted topics; 2025 and 2026 had none. The rationalised NCERT is the syllabus now.
\item Questions do come back (see the previous page), so the old papers are worth solving fully.
\end{itemize}\medskip
\textbf{\color{crimson}2. What to study from}\par\smallskip
\begin{enumerate}[leftmargin=1.3em,itemsep=2pt,topsep=0pt]
\item \textbf{NCERT Physics 2026-27} (XI Parts 1--2, XII Parts 1--2): the main book. Read every line, including the boxes, tables, figures, Summary and Points to Ponder.
\item \textbf{NCERT solved Examples and Exercises}: work all of them by hand. They are the single biggest source of numerical questions.
\item \textbf{NCERT Exemplar Physics}: for the applied questions; the MCQ and short-answer sections are enough.
\item \textbf{NCERT Lab Manual}: KCET 2022 and 2025 asked the vernier least count, the half-deflection formula for a galvanometer and real/apparent depth from it.
\item \textbf{KCET papers 2016--2026} (this book), solved under time, with the answer keys.
\item One standard reference for extra MCQ practice only after NCERT is complete. Do not start a new book in the last month.
\end{enumerate}\medskip
\textbf{\color{crimson}3. How to work}\par\smallskip
\begin{itemize}[leftmargin=1.1em,itemsep=2pt,topsep=0pt]
\item Make a one-page formula and fact sheet for each chapter while reading NCERT. Revise from these sheets.
\item Keep an error notebook: every wrong or guessed question, the reason, and the right idea. Read it every week.
\item Practise the graph, circuit and ray-diagram questions; about one question in four has a figure (153 of 660).
\item Learn units, dimensions, standard values ($e$, $h$, $c$, $m_e$, 13.6\,eV, $1\,\mathrm u = 931.5$\,MeV) and the ratio and proportion short cuts; they save time in nearly every paper.
\item Check every old question against the present syllabus; skip the ones marked \textsc{Deleted Portion}.
\end{itemize}
\columnbreak
\textbf{\color{crimson}4. Plan of action}\par\smallskip
{\small\renewcommand{\arraystretch}{1.25}
\begin{tabularx}{\linewidth}{@{}l>{\raggedright\arraybackslash}X@{}}\toprule
Phase & Work\\\midrule
1. Foundation & NCERT chapter by chapter, Class XII first (it carries about 70\% of the marks): text, Examples, Exercises, formula sheet.\\
2. Class XI & The same for Class XI, with more time on Motion in a Plane, Laws of Motion, Rotational Motion, Gravitation and Thermodynamics.\\
3. Topic drills & Solve this book topic by topic; start with the 15 most repeated topics. Then Exemplar MCQs on the same topics.\\
4. Full papers & One complete KCET paper every two days: 60 questions in 80 minutes. Check with the key and update the error notebook.\\
5. Last 2 weeks & Only revision: formula sheets, NCERT Summaries, error notebook, and a paper on alternate days. No new material.\\
\bottomrule\end{tabularx}}\par\medskip
\textbf{\color{crimson}5. In the examination hall}\par\smallskip
\begin{itemize}[leftmargin=1.1em,itemsep=2pt,topsep=0pt]
\item 80 minutes for 60 questions is 80 seconds each. First pass: answer every direct NCERT question at once (about half the paper).
\item Second pass: the numericals; leave any question that takes more than two minutes and come back.
\item KCET has had no negative marking, so mark an answer for every question before time is up.
\item Fill the OMR as you go, not at the end.
\end{itemize}\medskip
\textbf{\color{crimson}6. A weekly rhythm}\par\smallskip
{\small Five days of NCERT and practice, one day of a timed paper, one day of revision and the error notebook. A steady 2--3 hours of Physics a day, kept up for months, beats long sessions before the exam.}\par
\end{multicols}
